\documentclass[../../../include/open-logic-section]{subfiles}

\begin{document}

\olfileid[id]{his}{set}{mythology}

\olsection{Lebih Banyak Mitos daripada Sejarah?}

Ketika meninjau kembali peristiwa-peristiwa ini dengan jarak waktu lebih dari
satu abad, kita harus berhati-hati agar tidak menerima putusan-putusan
tersebut begitu saja. Hasil-hasilnya tentu baru, menarik, dan mengejutkan.
Namun, seberapa mengejutkankah hasil-hasil itu sebenarnya? Dan benarkah
mereka menunjukkan bahwa kita tidak seharusnya mengandalkan intuisi
geometris?

Mengenai persoalan keterkejutan, \citet{Gouvea2011} menunjukkan bahwa catatan
Cantor yang terkenal kepada Dedekind,
``\emph{je le vois, mais je ne le crois pas}'', diambil cukup jauh dari
konteksnya. Berikut bagian yang lebih lengkap dari konteks tersebut (dikutip
dari \citeauthor{Gouvea2011}):
\begin{quote}
Maafkan kegairahan saya terhadap pokok ini jika saya begitu banyak menuntut
kebaikan dan kesabaran Anda; komunikasi yang baru-baru ini saya kirimkan
kepada Anda begitu tidak terduga, begitu baru, bahkan bagi saya, sehingga
saya tidak dapat merasa tenang sebelum memperoleh dari Anda, sahabat yang
terhormat, suatu putusan mengenai kebenarannya. Selama Anda belum menyetujui
saya, saya hanya dapat berkata:
\emph{je le vois, mais je ne le crois pas.}
\end{quote}
Cantor mengetahui bahwa hasilnya ``begitu tidak terduga, begitu baru''.
Namun, diragukan bahwa ia pernah menganggap hasilnya \emph{mustahil
dipercaya}. Sebagaimana ditunjukkan \citeauthor{Gouvea2011}, ia hanya meminta
Dedekind memeriksa bukti yang telah diberikannya.

Mengenai persoalan intuisi geometris: Peano menerbitkan kurva pengisi
ruangnya tanpa menyertakan diagram apa pun. Namun, ketika Hilbert menerbitkan
kurvanya, ia menjelaskan tujuannya: ia hendak memberi pembaca suatu cara yang
jelas untuk memahami hasil Peano jika mereka ``menggunakan intuisi
geometris berikut''; lalu ia menyertakan serangkaian \emph{diagram} seperti
yang diberikan dalam \olref[his][set][pathology]{sec}.

Secara lebih umum, walaupun diagram agak ketinggalan mode dalam bukti yang
diterbitkan, tidak dapat disangkal bahwa para matematikawan \emph{sering}
menggunakan diagram ketika membuktikan sesuatu. (Secara kasar, matematikawan
yang baik mengetahui kapan mereka dapat mengandalkan intuisi geometris.)

Singkatnya: jangan percaya begitu saja pada kehebohannya; setidaknya, jangan
menerimanya tanpa pemeriksaan. Untuk pembahasan lebih lanjut, Anda dapat
membaca \citet{Giaquinto2007}.

\end{document}
